\chapter{Fix 30: The Talagrand T2 Inequality (Concentration)}
\label{ch:fix30_talagrand}

\section{Context: Non-Perturbative Tail Bounds}
A central difficulty in Constructive QFT for non-Abelian theories is controlling the "large field" region. In perturbation theory, fields are assumed to be small (Gaussian fluctuations). However, for validity of the global theory, we must prove that the probability of large fluctuations decays super-exponentially.
Standard Gaussian inequalities (Gross, hypercontractivity) often rely on convexity, which fails for the Yang-Mills action due to the commutator term $[A_\mu, A_\nu]^2$.
To resolve this, we utilize the modern framework of \textbf{Optimal Transport} to derive a \textbf{Talagrand Transport-Entropy Inequality (T2)}. This powerful inequality implies Gaussian concentration of measure even in the absence of global convexity, provided a curvature condition is met.

\section{Derivation: From Log-Sobolev to Transport}

\subsection{Step 1: The Wasserstein Distance}
We define the $L^2$-Wasserstein distance $W_2(\mu, \nu)$ between two probability measures $\mu, \nu$ on the space of lattice gauge fields $\mathcal{A}$:
\begin{equation}
W_2(\mu, \nu)^2 = \inf_{\pi \in \Pi(\mu, \nu)} \int_{\mathcal{A} \times \mathcal{A}} d_{geo}(A, A')^2 d\pi(A, A')
\end{equation}
where $d_{geo}$ is the geodesic distance on the group manifold $SU(N)^L$.

\subsection{Step 2: The Otto-Villani Theorem}
We have previously established the Log-Sobolev Inequality (LSI) with constant $\alpha$ (see Fix 2 and Section R.71). The Otto-Villani theorem states that for a Riemannian manifold equipped with a measure $e^{-S} dx$, LSI implies the Transport Entropy inequality $T_2(\alpha)$.
Specifically, if the measure $\mu$ satisfies LSI($\alpha$):
\begin{equation}
\text{Ent}_\mu(f^2) \le \frac{2}{\alpha} \int |\nabla f|^2 d\mu
\end{equation}
Then it satisfies the Talagrand Inequality:
\begin{equation}
W_2(\nu, \mu) \le \sqrt{\frac{2}{\alpha} H(\nu | \mu)}
\end{equation}
where $H(\nu|\mu)$ is the relative entropy (Kullback-Leibler divergence) of any probability measure $\nu$ with respect to the vacuum measure $\mu$.

\section{Implementation: Concentration of Measure}
The physical consequence of the $T_2$ inequality is the \textbf{Gaussian Concentration of Measure}.
Let $F: \mathcal{A} \to \mathbb{R}$ be any gauge-invariant observable that is strictly Lipschitz continuous with constant $L$ (e.g., a Wilson Loop $W_C$ or a flux tube operator).
By Marton's argument restricted to the transport setting:

\begin{theorem}[Non-Perturbative Concentration]
For the Yang-Mills vacuum measure $\mu$, and any $L$-Lipschitz observable $F$:
\begin{equation}
\mathbb{P}_\mu \left( |F - \mathbb{E}_\mu[F]| \ge t \right) \le 2 \exp\left( - \frac{\alpha t^2}{2 L^2} \right)
\end{equation}
\end{theorem}

\section{Implication for Stability}
This result is crucial for the "Hierarchical Zegarlinski" inductive step. It guarantees that the interaction terms between blocks, which act as effective disorders for the coarse-grained variables, have sub-Gaussian tails.
This prevents the "proliferation of large fields" that would otherwise destroy the analyticity of the renormalization group flow. It rigorously justifies the "Small Field Region" assumption used in Balaban's bounds, not by assumption, but as a theorem derived from the transport properties of the measure.


